\section{Direction Estimation and Patchwise Enhancement}
\label{sec:direction}

% PAGE TARGET: ~6 pages.
% Andrew-style narrative: a small number of displayed equations, supported
% by prose, with figures for every algorithm step.

\subsection{From a DRP Stack to a Per-Pixel Direction Map}
\label{sec:direction:perpixel}

The first stage of the pipeline reduces the four-dimensional DRP stack
to two two-dimensional maps that summarise, at each pixel, the strength
and direction of the local reflectance anisotropy. The motivation for
this reduction comes from a simple observation about laid lines. At a
pixel that sits on the peak of a laid line, the directional reflectance
profile shows a clear two-lobed pattern: the pixel returns more light
when illuminated from one side of the line normal than from the other,
and the magnitudes of the two lobes are well separated. At a pixel
between laid lines the same profile is much closer to isotropic, with
no preferred direction. Reducing each pixel's profile to a magnitude
of anisotropy and a dominant azimuth therefore captures most of the
information that a laid-line detector needs, and does so in a form
that is two orders of magnitude smaller than the raw stack.
% TODO: Figure 2.1 -- side-by-side DRP annular plots for a pixel on a
% laid-line peak and a pixel between laid lines, mirroring the
% structure of Andrew Figure 15. Caption: ``Pixel DRPs at a laid-line
% peak and a trough show qualitatively different anisotropy.''

The estimator chosen here is the first-order Fourier coefficient of
the $\phi$-marginalised reflectance, which is the simplest direction
estimator that respects the circular structure of the azimuth
coordinate. Concretely, the reflectance is first averaged over the
elevation angles to give a one-dimensional function of azimuth at
each pixel, and that function is then projected onto $\cos\phi$ and
$\sin\phi$:
\begin{equation}
  X(x,y) \;=\; \sum_{j=1}^{N_\phi}
      \bigl(p(x,y;\phi_j) - \bar p(x,y)\bigr)\cos\phi_j ,
  \qquad
  Y(x,y) \;=\; \sum_{j=1}^{N_\phi}
      \bigl(p(x,y;\phi_j) - \bar p(x,y)\bigr)\sin\phi_j ,
  \label{eq:dir_fourier}
\end{equation}
with $p(x,y;\phi)$ the elevation-averaged reflectance and $\bar p$ its
mean over $\phi$. The magnitude of anisotropy is then $M = \sqrt{X^2 + Y^2}$
and the dominant azimuth is $\hat\alpha = \mathrm{atan2}(Y, X)$. Both
maps are computed once per folio in a single vectorised pass over the
stack. The azimuth is retained in its full range
$\hat\alpha \in (-180^\circ, 180^\circ]$ rather than being folded
modulo $180^\circ$, because that is the native angular range of the
DRP acquisition and because the fold can be deferred to whichever
downstream stage actually needs it.

Two implementation details deserve mention. The first is that the
magnitude map is clipped at twice its own mean before being
normalised, which suppresses the small number of pixels (typically
sitting on ink or on adhesive tape) whose reflectance saturates at
all elevations and would otherwise dominate the dynamic range. The
second is that the elevation marginalisation in
\Cref{eq:dir_fourier} accepts an optional subsampling factor
$(s_\theta, s_\phi)$ that reduces the working size of the stack;
in the configuration used in this report the subsampling is set to
$(2, 2)$, which gives a four-fold reduction in compute cost with no
measurable degradation of the estimated azimuth, because the
Lambertian lobe is smooth on the scale of the angular sampling
lattice. An example output of the stage is shown in
\Cref{fig:direction_output} for one of the folios in the benchmark.
% TODO: Figure 2.2 -- a real folio image alongside its magnitude and
% azimuth maps. The azimuth map can use a circular colour scheme.

\subsection{Why Per-Pixel Azimuths Are Not Enough}
\label{sec:direction:noisy}

The per-pixel azimuth produced by the first stage is sufficient to
visualise the laid-line orientation by eye, but it is too noisy to
drive a frequency-domain detector directly. Two distinct sources
contribute to this noise. The first is photon noise in the source
images, which propagates into the first-order Fourier coefficients
of \Cref{eq:dir_fourier} and therefore into both the magnitude and
the azimuth. The second is the failure mode of the estimator itself
at low-anisotropy pixels: whenever the local Lambertian lobe is weak
relative to noise, $X$ and $Y$ are both small and the azimuth is
essentially uniform on $(-180^\circ, 180^\circ]$. A direct frequency
analysis of an image driven by such an azimuth field would inherit
the noise of these failure-mode pixels and would lose contrast
wherever the laid-line signal is weakest, which is precisely where
the detector most needs help.

The remedy adopted here is to aggregate the per-pixel azimuth over
larger spatial patches before mapping it to an enhanced greyscale
image. Each patch contributes a single best-guess direction, and
each pixel then receives a likelihood that its own azimuth agrees
with that of its enclosing patch. The output is a smooth greyscale
field in which the dynamic range is concentrated on regions whose
azimuth is consistent with their neighbourhood, and is suppressed
wherever the azimuth varies erratically.
% TODO: Figure 2.3 -- noisy per-pixel azimuth vs the smoother
% patchwise-aggregated greyscale, on the same folio, side by side.

\subsection{Patchwise Trigonometric Aggregation}
\label{sec:direction:patch}

The page is covered by overlapping square patches of side
$P = 512$ pixels, stepped at stride $S = 256$ in both the horizontal
and the vertical directions, so that each interior pixel falls inside
four patches. For each patch, the dominant direction is computed as
the circular mean of the per-pixel azimuths:
\begin{equation}
  \bar\alpha(\Omega) \;=\;
      \mathrm{atan2}\!\left(
        \sum_{(x,y)\in\Omega} \sin\hat\alpha(x,y),\;
        \sum_{(x,y)\in\Omega} \cos\hat\alpha(x,y)
      \right).
  \label{eq:circmean}
\end{equation}
The choice of circular distance rather than arithmetic distance is
essential because of the wraparound of the azimuth coordinate. A
patch in which half the pixels report $+179^\circ$ and the other
half report $-179^\circ$ has a true mean direction of $180^\circ$,
which is just one degree away from each individual pixel. The
arithmetic mean of those same values is $0^\circ$, which would point
at right angles to the actual laid lines. The trigonometric collapse
in \Cref{eq:circmean} avoids this trap by averaging the unit vectors
of the per-pixel azimuths rather than the azimuths themselves.

Each pixel in a patch is then assigned a Gaussian likelihood of
agreement with the patch direction, with angular tolerance
$\sigma = 10^\circ$ measured along the $360^\circ$ circle. The
per-patch likelihood fields are stitched back into a single image
using a two-dimensional Hann window as the blending weight, which
suppresses the discontinuities that would otherwise appear at patch
boundaries. The overlap factor $P/S = 2$ used in this report is
sufficient in practice for the folios in the benchmark, in which
the dominant direction varies by less than five degrees across any
single patch. A final local contrast enhancement step (a small
Gaussian smoothing followed by contrast-limited adaptive histogram
equalisation) conditions the resulting greyscale field for the
frequency-domain detector of \Cref{sec:detection}. An example of
the enhanced image alongside the original azimuth field is shown
in \Cref{fig:patchwise_output}.
% TODO: Figure 2.4 -- patchwise direction field as a vector overlay,
% plus the enhanced greyscale image and its histogram before/after
% CLAHE. Cf. Andrew Figure 24 (vector-field refinement).

The output of this stage is the field $L(x, y) \in [0, 1]$, treated
as an eight-bit greyscale image for the detection stage that follows.
The patch direction field $\bar\alpha(\Omega)$ is retained as a
by-product, and is reused both by the visualisation overlays in
\Cref{sec:benchmark} and by the automatic line-direction detector
of \Cref{sec:detection:auto}.
